% 狮子座（综述）
% license CCBYSA3
% type Wiki

本文根据 CC-BY-SA 协议转载翻译自维基百科\href{https://en.wikipedia.org/wiki/Leo_(constellation)}{相关文章}。

狮子座（Leo，/ˈliːoʊ/）是黄道十二星座之一，位于西侧的巨蟹座（Cancer，螃蟹）与东侧的室女座（Virgo，处女）之间，坐落在北天球。其名称源自拉丁语，意为“狮子”，在古希腊文化中代表神话英雄赫拉克勒斯（Heracles）在其十二项功绩之一中所杀死的尼米亚狮子。其古老的天文符号为♌︎。作为公元2世纪天文学家托勒密所描述的48个星座之一，狮子座至今仍是现代88个星座之一，并且由于其众多明亮的恒星以及一个极具辨识度、令人联想到伏地狮子的独特形状，狮子座也是最容易识别的星座之一。
\subsection{特征}
\subsubsection{恒星}
\begin{figure}[ht]
\centering
\includegraphics[width=6cm]{./figures/cd1a625b724681d2.png}
\caption{狮子座的肉眼可见外观（图像中央的明亮天体为2004年3月的行星木星）。} \label{fig_Leo_1}
\end{figure}
狮子座包含许多明亮恒星，其中许多在古代就已被单独识别。该星座中有九颗明亮恒星可以轻易用肉眼观测到，其中四颗属于一等或二等星，这使得狮子座显得尤为醒目。这九颗恒星中的六颗还组成了一个称为“镰刀”（The Sickle）的星群，对于现代观测者来说，其形状类似一个反向的“问号”。镰刀由六颗恒星标记：狮子座$\varepsilon$、狮子座$\mu$、狮子座$\zeta$、狮子座$\gamma$、狮子座$\eta$以及狮子座$\alpha$。其余三颗恒星构成一个等腰三角形，其中狮子座$\beta$（Denebola）标示狮子的尾巴，而狮子的其余身体部分则由狮子座$\delta$和狮子座$\theta$勾勒出来$^{[2]}$。
\begin{itemize}
\item 狮子座$\alpha$，称为轩辕十四（Regulus），是一颗蓝白色主序星，视星等为1.34，距离地球77.5光年。它是一对双星，用双筒望远镜即可分辨，其伴星的视星等为7.7。其传统名称Regulus意为“小国王”。
\item 狮子座$\beta$，又称Denebola，位于星座中与轩辕十四相对的一端。它是一颗蓝白色恒星，视星等为2.23，距离地球36光年。Denebola一名意为“狮子的尾巴”。
\item 狮子座$\gamma$，又称Algieba，是一个双星系统，并伴有一颗视线上的第三成员；主星和伴星可用小型望远镜分辨，第三成员可用双筒望远镜看到。主星是一颗金黄色巨星，视星等为2.61，伴星性质相似但视星等为3.6；它们的轨道周期约为600年，距离地球126光年。无关的第三星，即狮子座40，是一颗略带黄色的恒星，视星等为4.8。其传统名称Algieba意为“前额”。
\item 狮子座$\delta$，称为Zosma，是一颗蓝白色恒星，视星等为2.58，距离地球58光年。
\item 狮子座$\mu$，称为Rasalas，是一颗红色恒星，视星等为3.88，亮度足以用肉眼观测。该系统距离地球124光年。Rasalas亦称Alshemali，这两个名称均源自“Al Ras al Asad al Shamaliyy”的缩写，意为“狮子朝南的头部”。
\item 狮子座$\theta$，称为Chertan，是一颗白色恒星，视星等为3.324，可用肉眼观测，是该星座中较亮的恒星之一，距离地球约165光年。Chertan一名源自阿拉伯语al-kharātān，意为“两根小肋骨”，最初指的是狮子座$\delta$和狮子座$\theta$，$^{[3]}$。
\end{itemize}
\begin{figure}[ht]
\centering
\includegraphics[width=8cm]{./figures/a1a3d9c69259c7b9.png}
\caption{狮子座最亮的恒星} \label{fig_Leo_2}
\end{figure}
\begin{figure}[ht]
\centering
\includegraphics[width=6cm]{./figures/fe6607081f71245d.png}
\caption{狮子座，上方是小狮子座，如乌拉尼亚之镜所描绘，这是一套约1825年在伦敦出版的星座卡片} \label{fig_Leo_3}
\end{figure}
狮子座中还包含一颗明亮的变星——红巨星狮子座R（R Leonis）。它是一颗米拉型变星，最暗时视星等约为10，通常的最亮视星等为6；并且会周期性地增亮至视星等4.4。狮子座R距离地球约330光年，周期为310天，直径约为太阳直径的450倍$^{[4]}$。

恒星沃尔夫359（Wolf 359，又名狮子座CN）位于狮子座中，是距离地球最近的恒星之一，距离仅7.8光年。沃尔夫359是一颗红矮星，视星等为13.5；由于它是一颗耀斑星，其亮度会周期性地增亮不超过一个星等$^{[4]}$。格利泽436（Gliese 436）是一颗位于狮子座内、距离太阳约33光年的暗弱恒星，其周围运行着一颗凌日的海王星质量级系外行星$^{[5]}$。

碳星CW Leo（IRC +10216）是在红外N波段（波长10\,$\mu$m）夜空中最明亮的恒星。

恒星SDSS J102915+172927（又称Caffau之星）是一颗位于狮子座方向的银河系晕族II恒星，其年龄约为130亿年，是银河系中已知最古老的恒星之一。它具有目前已知恒星中最低的金属丰度。

现代天文学家（包括1602年的第谷·布拉赫）将一组原本构成“狮子尾端毛簇”的恒星从狮子座中剔除，并用它们建立了新的星座——后发座（Berenice之发），尽管在古希腊和古罗马时期，这一名称已有先例$^{[6]}$。
\subsubsection{深空天体}
狮子座中包含许多明亮的星系，其中最著名的包括梅西耶65、梅西耶66、梅西耶95、梅西耶96、梅西耶105以及NGC\,3628$^{[7]}$，其中前两者属于著名的“狮子座三重星系”$^{[8]}$。

狮子座环（Leo Ring）是一团由氢和氦气体组成的云，位于该星座内两座星系的轨道附近。
\begin{figure}[ht]
\centering
\includegraphics[width=6cm]{./figures/3b150d76df662684.png}
\caption{} \label{fig_Leo_4}
\end{figure}
M66 是一座螺旋星系，属于狮子座三重星系，其另外两个成员是 M65 和 NGC\,3628。它距离地球约 3700 万光年，由于与三重星系中其他成员之间的引力相互作用，其形态略显扭曲，这些相互作用正在将恒星从 M66 中拉离。最终，最外层的恒星可能会形成一座绕 M66 运行的矮星系$^{[9]}$。M65 和 M66 均可通过大型双筒望远镜或小型望远镜观测到，但它们高度集中的核区以及拉伸的形态特征，只有在大型业余观测设备中才能清晰辨认$^{[4]}$。
\begin{figure}[ht]
\centering
\includegraphics[width=6cm]{./figures/04a0b3e20a43eeff.png}
\caption{} \label{fig_Leo_5}
\end{figure}